\section{Amplification to rank two}\label{sec:amplification}
We recall the amplification convention for a $\mathrm{II}_1$ factor
$(Q,\tau_Q)$. Let $\Tr$ be the usual semifinite trace on $B(\ell^2)$,
normalized to give every rank-one projection trace one. For $t>0$,
choose a projection $p$ in the stabilized factor with
$(\tau_Q\otimes\Tr)(p)=t$ and set
\[
 Q^t=p\bigl(Q\,\overline\otimes\,B(\ell^2)\bigr)p,
 \qquad
 \tau_{Q^t}=t^{-1}(\tau_Q\otimes\Tr)\big|_{Q^t}.
\]
Projections of every positive finite trace exist in this stabilized
factor. Equivalence of projections of equal finite trace makes the
normalized corner independent of the choice of $p$, up to a normal
trace-preserving $*$-isomorphism.

Amplification respects normal tracial isomorphisms. Indeed, if
$\Theta:Q\to R$ is such an isomorphism, the normal tensor-product map
$\Theta\otimes\mathrm{id}$ preserves the semifinite trace and sends
$p$ to a projection $p'$ of the same trace $t$. Its restriction is an
onto normal isomorphism between the corresponding corners. It sends
the corner unit $p$ to $p'$ and preserves their traces normalized by
$t^{-1}$. We denote the resulting map by $\Theta^t:Q^t\to R^t$.

For real $s>1$, let $N_s=L(\F_s)$ be the interpolated free group factor,
with the usual group factor at integer parameters. The established
amplification formula is
\begin{equation}\label{eq:interpolation-formulas}
 N_s^t\cong N_{1+(s-1)t^{-2}}
 \qquad(s>1,\ t>0);
\end{equation}
see \cite[Theorem~2.4]{dykema1994}. This is an unconditional
isomorphism formula for the interpolated family and does not assume
that distinct parameters determine distinct isomorphism classes.
It is the only interpolation theorem used in the proof of
Theorem~\ref{thm:main}.

\begin{proof}[Proof of Theorem~\ref{thm:main}]
Apply Theorem~\ref{thm:rank-step} with $n=3$ and $n=4$ to obtain
\[
 N_3\cong N_4\cong N_5.
\]
Let $\Theta:N_3\to N_5$ be the composition of these normal
trace-preserving isomorphisms. Amplifying by $\sqrt2$ gives a
unital normal trace-preserving isomorphism
\[
 \Theta^{\sqrt2}:N_3^{\sqrt2}\longrightarrow N_5^{\sqrt2}.
\]
Formula~\eqref{eq:interpolation-formulas} identifies its domain and
range as follows:
\[
 N_3^{\sqrt2}\cong N_{1+(3-1)/2}=N_2,
 \qquad
 N_5^{\sqrt2}\cong N_{1+(5-1)/2}=N_3.
\]
These identifications and the amplified map compose to give
$\Phi:L(\F_2)\to L(\F_3)$. Every map in the composition is a unital
normal trace-preserving $*$-isomorphism. Thus $\Phi$ is bijective
and has all the asserted properties.
\end{proof}


For completeness, we record the consequences for the whole interpolated
family. Put $N_\infty=L(\F_\infty)$ and, for a $\mathrm{II}_1$ factor $Q$,
define its fundamental group by
\[
 \mathcal F(Q)=\{t>0:Q^t\cong Q\}.
\]
Here and below the isomorphisms are normal and preserve the normalized
traces. The classical free group factor alternative says that either
all $N_r$, $1<r\leq\infty$, are isomorphic, or $N_r\not\cong N_s$
whenever $1<r<s\leq\infty$. The inclusion of the infinite parameter
uses R\u adulescu's dichotomy \cite[Corollary~4.7]{radulescu1994},
restated explicitly in \cite[Theorem~1.1]{goldbring_pi2025}.
Dykema's corresponding dichotomy \cite[Corollary~4.2]{dykema1994}
is stated for finite parameters.

\begin{corollary}\label{cor:all-interpolated}
For all $1<r,s\leq\infty$ there is a unital normal trace-preserving
$*$-isomorphism $L(\F_r)\cong L(\F_s)$. Moreover,
\[
 \mathcal F\bigl(L(\F_r)\bigr)=\R_{>0}
 \qquad(1<r\leq\infty).
\]
\end{corollary}
\begin{proof}
Theorem~\ref{thm:main} rules out the pairwise-nonisomorphic alternative,
so all parameters, including $\infty$, give the same isomorphism class.
For finite $r>1$ and every $t>0$, formula~\eqref{eq:interpolation-formulas}
then gives
\[
 N_r^t\cong N_{1+(r-1)t^{-2}}\cong N_r.
\]
Thus $\mathcal F(N_r)=\R_{>0}$. Finally, $N_\infty\cong N_2$ and
functoriality of amplification give
$N_\infty^t\cong N_2^t\cong N_2\cong N_\infty$ for every $t>0$.
\end{proof}
